
\medskip
\phantomsection\label{ctr:sga3-VIB-1.2.2}%
\noindent\textbf{Counter\SGARefLink{ctr:sga3-VIB-1.2.2}{example 1.2.2}.---}
\emph{Let \(k\) be a field of characteristic zero, let \(G\) be
the constant \(k\)-group \(\mathbb Z\), and let \(H\) be the
additive \(k\)-group \(\mathbb G_{a,k}\). Let \(u:G\to H\) be a
morphism of \(k\)-groups. If \(u\ne0\), then \(u(G)\) is not closed
in \(H\).}

\medskip
\phantomsection\label{prop:sga3-VIB-1.3}%
\noindent\textbf{Proposition 1.3.---}%
\footnote{\label{note:sga3-VIB-5}Editor's note: in the original,
\SGARefLink{prop:sga3-VIB-1.3}{1.3}, \SGARefLink{cor:sga3-VIB-1.3.1}{1.3.1}, and \SGARefLink{lem:sga3-VIB-1.3.1.1}{1.3.1.1} were stated with \(A\) a field. For
completeness, we treat the case of an Artinian local ring, and for
this purpose have added \hyperref[lem:sga3-VIB-1.3.0]{Lemma~1.3.0}.}
\emph{Let \(A\) be an Artinian local ring with residue field \(k\),
let \(G\) be an \(A\)-group locally of finite type and flat, and
let \(u:G\to H\) be a morphism of \(A\)-groups. The following
assertions are equivalent:
\begin{enumerate}
\item[\textup{(i)}]\phantomsection\label{concl:sga3-VIB-1.3-i} \(u\) is flat (respectively, quasi-finite,
  unramified, smooth, or \'etale) at some point of \(G\);%
  \footnote{\label{note:sga3-VIB-6}Editor's note: recall that a
  morphism of schemes \(f:X\to Y\) is said to be of finite type
  (respectively, of finite presentation or quasi-finite) at \(x\)
  if there are affine open subsets \(V\), containing \(f(x)\), and
  \(U\), containing \(x\), such that \(f(U)\subset V\) and
  \(\mathcal O_X(U)\) is an \(\mathcal O_Y(V)\)-algebra of finite
  type (respectively, of finite presentation, or with the
  additional condition that \(f^{-1}(f(x'))\) is finite for every
  \(x'\in U\)); see also Editor's note 40 of Expos\'e~V. The
  morphism is locally quasi-finite (respectively, locally of finite
  type or locally of finite presentation) if it has the property
  at every point \(x\). It is smooth (respectively, unramified or
  \'etale) at \(x\) if there is an open neighborhood \(U\) of \(x\)
  such that \(f|_U:U\to Y\) is smooth (respectively, unramified or
  \'etale). These definitions show that the loci where \(f\) is of
  finite presentation, of finite type, quasi-finite, smooth,
  unramified, or \'etale are open in \(X\). Moreover, because the
  flatness locus of a morphism locally of finite presentation is
  open (EGA~IV\(_3\), 11.3.1), the locus where \(f\) is both of
  finite presentation and flat is open as well. These facts will
  be used repeatedly below.}
\item[\textup{\SGAInlineRefTarget{prop:sga3-VIB-1.3:item:ii}(ii)}] \(u\) is flat (respectively, quasi-finite,
  unramified, smooth, or \'etale).
\end{enumerate}}

\emph{Proof.}
It is enough to prove that \textup{\SGARefLink{concl:sga3-VIB-1.3-i}{(i)}} implies \textup{\SGARefLink{prop:sga3-VIB-1.3:item:ii}{(ii)}}. We
first establish the following lemma.

\medskip
\phantomsection\label{lem:sga3-VIB-1.3.0}%
\noindent\textbf{Lemma 1.3.0.---}
\emph{Let \(A\to B\to C\) be morphisms of commutative rings, and
let \(\mathfrak n\) be a nilpotent ideal of \(A\). Suppose that
\(C/\mathfrak nC\) is a
\((B/\mathfrak nB)\)-algebra of finite type.
\begin{enumerate}
\item[\textup{\SGAInlineRefTarget{lem:sga3-VIB-1.3.0:item:i}(i)}] Then \(C\) is a \(B\)-algebra of finite type.
\item[\textup{\SGAInlineRefTarget{lem:sga3-VIB-1.3.0:item:ii}(ii)}] Moreover, if \(C\) is flat over \(A\) and
  \(C/\mathfrak nC\) is a
  \((B/\mathfrak nB)\)-algebra of finite presentation, then \(C\)
  is a \(B\)-algebra of finite presentation.
\end{enumerate}}

Indeed, let \(x_1,\ldots,x_n\) be elements of \(C\) whose images
generate \(C/\mathfrak nC\) as a
\((B/\mathfrak nB)\)-algebra. By the nilpotent Nakayama lemma, the
\(x_i\) generate \(C\) as a \(B\)-algebra. This proves \SGARefLink{lem:sga3-VIB-1.3.0:item:i}{(i)}. Let
\[
  \varphi:B[X_1,\ldots,X_n]\longrightarrow C
\]
be the resulting surjective morphism, and put \(I=\ker\varphi\).

Suppose now that \(C\) is flat over \(A\) and that
\(C/\mathfrak nC\) is of finite presentation over
\(B/\mathfrak nB\). On the one hand, \(I/\mathfrak nI\) identifies
with the kernel of
\[
  \overline\varphi
    =\varphi\otimes_A(A/\mathfrak n).
\]
On the other hand, by EGA~IV\(_1\), 1.4.4,
\(\ker\overline\varphi\) is an ideal of finite type. Choose
\(P_1,\ldots,P_s\in I\) whose images generate
\(I/\mathfrak nI\) as an ideal. The nilpotent Nakayama lemma shows
that the \(P_i\) generate \(I\). This proves \SGARefLink{lem:sga3-VIB-1.3.0:item:ii}{(ii)}.

We return to the proof of \hyperref[prop:sga3-VIB-1.3]{1.3}. Let \(x\) be an arbitrary point of
\(G\). Since \(G\) is flat over \(A\), the fiberwise flatness
criterion, in the form EGA~IV\(_3\), 11.3.10.2, shows that \(u\)
is flat at \(x\) if \(u\otimes_Ak\) is. Likewise, the preceding
lemma shows that \(u\) is of finite type (respectively, of finite
presentation) at \(x\) if \(u\otimes_Ak\) is. The remaining
properties can then be checked on the fibers
(cf.\ EGA~IV\(_4\), 17.4.1, 17.5.1, and 17.6.1 for the
unramified, smooth, and \'etale cases). We are therefore reduced to
the case \(A=k\).

Now let \(x\) be a point of \(G\) at which one of the conditions
in \hyperref[concl:sga3-VIB-1.3-i]{1.3(i)} holds. The properties under consideration are preserved
by \(\mathrm{fpqc}\) descent (cf.\ EGA~IV, 2.5.1, 2.7.1, and
17.7.1). Thus, after replacing \(k\) by an algebraic closure of
\(\kappa(x)\), we reduce to the case in which \(k\) is
algebraically closed and \(x\in G(k)\).

Since \(G\) is a Jacobson scheme (cf.\ EGA~IV\(_3\), 10.4.7), and
the set \(W\) of points of \(G\) at which \(u\) is flat
(respectively, quasi-finite, unramified, smooth, or \'etale) is
stable under generization%
\footnote{\label{note:sga3-VIB-7}Editor's note: the original stated
that the flat locus \(W\) is open, which does not seem evident
a priori.}
(respectively, open), it is enough to show that every
\(y\in G(k)\) belongs to \(W\). For such a point \(y\), the
translation
\[
  r_y\circ r_x^{-1}
\]
sends \(x\) to \(y\), so \(u\) has the required property at \(y\).
Thus \(y\in W\), and the proof is complete.
